\documentclass[11pt,a4paper]{article}
% Common preamble for the hint sheets (\input from each file)
\usepackage[utf8]{inputenc}
\usepackage[T1]{fontenc}
\usepackage{lmodern}
\usepackage{amsmath,amssymb,amsthm,mathtools}
\usepackage[a4paper,margin=2.3cm]{geometry}
\usepackage{enumitem}
\usepackage{xcolor}
\usepackage[hidelinks]{hyperref}
\renewcommand{\leq}{\leqslant}
\renewcommand{\le}{\leqslant}
\renewcommand{\geq}{\geqslant}
\renewcommand{\ge}{\geqslant}
\setlength{\parindent}{0pt}
\setlength{\parskip}{0.35em}
% \hintstitle{Sheet title}{Class work / Homework}
\newcommand{\hintstitle}[2]{\begin{center}{\LARGE\bfseries Hints: #1\par}\vspace{0.5em}{\large #2\par}\end{center}
\noindent\textit{The hints to each problem are cumulative. Read one, try again, and open the next only if you are still stuck.}\par\medskip\hrule\medskip}
\newcommand{\problem}[1]{\par\bigskip{\large\bfseries Problem #1}\par\nopagebreak\smallskip}
\newcommand{\subpart}[1]{\par\smallskip\textbf{(#1)}\par\nopagebreak}
\newcommand{\hint}[1]{\par\smallskip\noindent\textcolor{black!70}{\textbf{Hint #1.}}\enspace}
\newcommand{\note}[1]{\par\smallskip\noindent\textit{Note.} #1\par}
\providecommand{\lt}{<}
\providecommand{\gt}{>}
\providecommand{\tightlist}{\setlength{\itemsep}{0pt}\setlength{\parskip}{0pt}}
\emergencystretch=3em

\begin{document}
\hintstitle{Number theory -- 2}{Class work}

\problem{1}
\hint{1} Theorem 1: write $m=q\,\mathrm{ord}_n(a)+r$ with $0\le r<\mathrm{ord}_n(a)$; then $a^r\equiv1$, so $r=0$. If $p\mid 2^p-1$ for a prime $p$, then $\mathrm{ord}_q(2)$ (for a prime $q\mid 2^p-1$) divides $p$ and is not $1$, so it equals $p$. Hence $p\mid q-1$, and $q$ is odd.
\hint{2} Let $q\mid F_k=2^{2^k}+1$. Then $2^{2^{k+1}}\equiv1$ but $2^{2^k}\equiv-1$, so $\mathrm{ord}_q(2)=2^{k+1}$. To gain one more factor of $2$, find a square root of $2$ mod $q$: for $k\ge2$, check that $s=2^{2^{k-2}}(2^{2^{k-1}}-1)$ satisfies $s^2\equiv2\pmod q$.
\hint{3} Then $\mathrm{ord}_q(s)=2^{k+2}$, which divides $q-1$. For primes $\equiv1\pmod n$: if $q\mid\Phi_n(a)$ and $q\nmid n$, then $\mathrm{ord}_q(a)=n$, so $q\equiv1\pmod n$. Apply this to $a=N!\,n$, where $\Phi_n(a)\equiv\Phi_n(0)=\pm1\pmod{N!\,n}$, so all its prime factors exceed $N$.

\problem{2}
\hint{1} For $d\mid p-1$ let $\psi(d)$ be the number of elements of order $d$. If $\psi(d)>0$, these elements are among the $d$ roots of $x^d=1$, which form $\langle g\rangle$ for any one of them. So $\psi(d)\le\varphi(d)$, and $\sum_d\psi(d)=p-1=\sum_d\varphi(d)$ forces equality.
\hint{2} Prime powers: if $g$ is a primitive root mod $p$, then $g$ or $g+p$ has order $p(p-1)$ mod $p^2$, and such an element stays primitive mod every $p^k$ (use LTE on $g^{p-1}-1$). For $2p^k$ take an odd primitive root mod $p^k$. For other $n$, CRT gives a product of two groups of even order, or $(\mathbb Z/2^k)^*$ with $k\ge3$, and these are not cyclic.
\hint{3} $\mathrm{ord}_9(2)=6$, and $v_3(2^{2\cdot3^j}-1)=v_3(4-1)+j=1+j$, so $\mathrm{ord}_{3^k}(2)=2\cdot3^{k-1}=\varphi(3^k)$. For $k\ge2$, $2^j-1<2^k-1$ when $0<j<k$, so the order of $2$ modulo $2^k-1$ is exactly $k$.

\problem{3}
\hint{1} Make the partial sums telescope: try to choose $a_i$ so that $a_ia_{i+1}=c_i-c_{i+1}$ for a simple injective sequence $c_i$.
\hint{2} Take $a_i\equiv i^{-1}\pmod p$. Then $a_ia_{i+1}=\frac1i-\frac1{i+1}$, so the partial sums are $1-\frac1{j+1}$ for $j\le p-2$. These are distinct and avoid $0$ and $1$.
\hint{3} The last sum adds $a_{p-1}a_1=-1$ to $S_{p-2}=2$, giving $1$, which is new. Wonkiness: $a_ia_{i+2}-a_{i+1}^2=\frac{1}{i(i+1)^2(i+2)}\neq0$, and the two wrap-around values equal $-\frac32\ne0$ for $p>3$.

\problem{4}
\hint{1} The right-hand side counts ordered bases of $\mathbb F_p^n$: choose the $i$-th column outside the span of the previous ones, in $p^n-p^{i-1}$ ways. So it is $|\mathrm{GL}_n(\mathbb F_p)|$.
\hint{2} Find a subgroup of order $(p-1)^nn!$: the monomial matrices (permutation matrices times invertible diagonal matrices). Then apply Lagrange's theorem.

\problem{5}
\hint{1} The $b_n$ form an arithmetic progression $b_0+nr$ of length $p-4$, so $1,a,\dots,a^{p-5}$ are distinct, and $\mathrm{ord}_p(a)\ge p-4$. For $p>7$ this forces $\mathrm{ord}_p(a)=p-1$.
\hint{2} The complement of an AP of length $p-4$ in $\mathbb F_p$ is an AP of length $4$ with the same difference. The complement of $\{1,a,\dots,a^{p-5}\}$ is $\{0,u,u^2,u^3\}$ with $u=a^{-1}$. Up to reversal it equals $\{0,d,2d,3d\}$ or $\{-d,0,d,2d\}$.
\hint{3} Write $u=\alpha d$, $u^2=\beta d$, $u^3=\gamma d$. Then $\beta^2\equiv\alpha\gamma$, where $\{\alpha,\beta,\gamma\}=\{1,2,3\}$ or $\{-1,1,2\}$. Check the few cases: only $9\equiv2$, i.e.\ $p=7$, survives with $p>5$. For $p=7$, $a=3$ or $5$ works. Answer: $p=7$.
\end{document}
